\documentclass[runningheads]{llncs}
\usepackage[T1]{fontenc}
\usepackage{amsmath,amssymb}
\usepackage{booktabs,array}
\usepackage{xcolor}
\usepackage{listings}
\usepackage{xurl}
\usepackage[hidelinks]{hyperref}

\newcommand{\guardcert}{GuardCert}
\newcommand{\code}[1]{{\small\nolinkurl{#1}}}
\newcommand{\pending}[1]{\par\noindent\textcolor{red!65!black}{\textbf{Pending evidence.} #1}\par}
\newcommand{\run}[3]{#1\mathrel{\Downarrow}_{#2}#3}
\newcommand{\checkrun}[3]{#1\mathrel{\Downarrow}_{#2}#3}
\lstset{basicstyle=\ttfamily\small,columns=fullflexible,keepspaces=true,
  showstringspaces=false,frame=single,breaklines=true}

\begin{document}
\title{GuardCert: Verified Optimistic Transformation with Safe Runtime Conditions}
\titlerunning{Verified Optimistic Transformation}
\author{Anonymous Authors}
\authorrunning{Anonymous Authors}
\institute{Anonymous Institution}
\maketitle

\begin{abstract}
An optimistic transformation replaces a program fragment when a runtime
condition establishes the assumptions of a conditional correctness proof.
The replacement must preserve behavior even when the condition refuses.
Constructing that condition raises a separate verification problem: its reads,
arithmetic, and private state must be safe before the assumptions have been
established. We present \guardcert, a language-independent certificate interface
for local guarded transformations, together with a CompCert/Clight instance.
The interface separates check execution, conditional transformation proofs,
and the language-specific installation of replacements into complete programs.
The current implementation connects supported affine loop candidates and
loaded-root bounds to a complete-program simulation theorem. Nested loaded bounds
motivate additional services for ordered capture and source-prefix reasoning;
their complete optimizer integration remains in progress.
\keywords{Verified compilation \and Runtime conditions \and Polyhedral transformation}
\end{abstract}

\noindent\textbf{Working draft, 2026-10-06.}
This manuscript records proved interfaces and established execution paths.
Red paragraphs identify results still required for the intended case study.
This is not a submission-ready claim of complete OLO coverage or speedup.

\section{Introduction}
\label{sec:introduction}

A loop transformation may require stable bounds, non-overlapping memory
accesses, or arithmetic that agrees with mathematical integers. Static analysis
may fail to establish these facts for a useful candidate. A compiler can retain
the original fragment as a fallback and execute the candidate only after a
runtime check accepts. This choice also supports a candidate proposed by a
programmer or an external tool: the producer need not be the compiler pass that
eventually installs it.

The motivating object is a source fragment $S$, a candidate $T$, and a semantic
premise $A$ under which $T$ preserves the relevant behavior of $S$. An executable
condition $G$ establishes a sufficient entry property $B$, with $B\Rightarrow A$.
The transformed fragment selects $T$ when $G$ accepts and $S$ otherwise. This
composition rule is simple. The difficult question is how to obtain reusable
certificates for its premises on an actual language and program.

In particular, checking an assumption can require the assumption itself.
Consider an inner bound loaded from memory. Hoisting that load requires evidence
that the source reaches it. Treating its captured value as constant requires
evidence that intervening stores preserve the observed location. Checking those
stores may require addresses whose validity follows only from the source's
actual execution. A proof that starts with the entire cached loop execution
would therefore assume a result the guard is meant to establish.

Optimistic Loop Optimization (OLO) already addresses assumption construction,
condition simplification, and safe conditional preloading in a production
polyhedral setting~\cite{olo2017}. Our purpose is to investigate the mechanical
proof interfaces connecting these obligations, conditional candidate validation,
and complete-program compilation. The CompCert case study supplies concrete
machine integers, memory permissions, and control semantics~\cite{compcert2009}.

\paragraph{Research question.}
Can a small language-independent interface support a substantial guarded loop
transformation while reusable language and domain services discharge the
condition's safety, sufficiency, and installation obligations? The intended
contribution lies in this verified chain and its reuse. A generic conditional
composition theorem, insertion of an \code{if}, or use of the C layer alone does
not establish novelty relative to existing verified transformation systems.

\paragraph{Design and current results.}
\guardcert\ divides responsibility among three parties. The generic kernel
combines certificates for checks and local transformations. A language/IR
instance proves concrete execution, guarded choice, private-state transport,
and legal installation. A transformation implementer selects fragments and
supplies candidate proofs, source/model correspondence, and domain-specific
assumption derivation. In the present CompCert instance, supported affine
candidates, dependency checks, runtime conditions, source fallback, and backend
compilation are connected to whole-program simulation theorems. For a nested
loaded-bound profile, ordered capture and source-prefix advancement produce
the model and private entry relation needed by the guarded candidate proof.
The extracted compiler now installs interchange and tiling on the adapted
three-level OLO source, with array outputs, public exits, and executed store
order checked separately. Additional array bodies and control contexts use the
same compiler, including actual dependence refusal and repeated sites.

Condition reuse has an explicit practical limit. For a two-loaded-bound,
two-store successor, verified repeated-probe elimination preserves the original
condition's answer and reduces setup tests. Complete-call measurements on
selected warm inputs nevertheless remain slower than source while the
existing scans retain their work. The evaluation reports these negative results;
the next requirement is a cheaper safe sufficient condition, not further code
size reduction alone.

\pending{Establish a fresh-build path, extend and state the final supported
source class, condition algorithms, and evidence for proof reuse. Measure
broader condition cost and useful acceptance separately from code size before claiming
practical benefit.}

\paragraph{Scope.}
The current concrete target is sequential, behavior-preserving transformation
with entry fallback. The kernel abstracts over states, observations, and their
relation, but each host must prove the laws for the behaviors it exposes.
Checkpoint deoptimization, concurrent interference, approximate computation,
and patch specifications require different host or correctness contracts; they
are not established by the present Clight instance.

\section{A Small Certificate Kernel}
\label{sec:framework}

This section describes the implemented interfaces in
\code{prototype/interface/GuardInterface.v}. State, code, checks, and
observations are parameters. The kernel neither interprets C expressions nor
extracts assumptions from arbitrary source/candidate pairs.

\subsection{Host Semantics and Private Check State}

A \code{guard_host} supplies relations $\run{c}{s}{o}$ for code execution and
$\checkrun{g}{s}{(b,s')}$ for check execution, together with a safety predicate
and a guarded-choice constructor $\operatorname{select}(g,t,f)$. Its dispatch law
is exact:
\begin{equation}
\run{\operatorname{select}(g,t,f)}{s}{o}
\ \Longleftrightarrow\
\exists b,s'.\ \checkrun{g}{s}{(b,s')}\ \land\
\run{(\mathbf{if}\ b\ \mathbf{then}\ t\ \mathbf{else}\ f)}{s'}{o}.
\label{eq:dispatch}
\end{equation}
The language instance proves this law for its declared observations. It must
account for observable check effects and any exceptional or divergent behavior
it claims to include. The constructor need not be a single syntactic conditional.

Checks may alter private temporaries. Consequently, the interface retains both
the original entry $s$ and the checked state $s'$. Accepted and refused entry
relations $R_+(s,s')$ and $R_-(s,s')$ state what the selected branch may rely on.
For a read-only frontend, these relations can be equality. For a private cursor
scan, they instead express preservation of public temporaries and memory.

\subsection{Guard and Transformation Certificates}

A \code{guard_certificate} is parameterized by an entry domain $D$, premise $P$,
the two entry relations, and the actual check $g$. It requires:
\begin{align}
D(s)&\Rightarrow\operatorname{safe}(g,s),\label{eq:safety}\\
D(s)&\Rightarrow\exists b,s'.\ \checkrun{g}{s}{(b,s')},\label{eq:available}\\
D(s)\land\checkrun{g}{s}{(b,s')}&\Rightarrow
\begin{cases}P(s)\land R_+(s,s')&b=\mathit{true},\\R_-(s,s')&b=\mathit{false}.
\end{cases}\label{eq:sound}
\end{align}
The premise is anchored at the original entry. Refusal establishes a valid
fallback entry, not $\neg P$. A conservative condition is therefore permitted.
Availability is existential; it is not a general check termination theorem.
Safety is a host predicate whose intended interpretation the language must
justify separately.

A \code{conditional_certificate} proves target-to-source refinement for both
branches. Let $o_t\preceq o_s$ be the supplied observation relation. Its accepted
clause is
\begin{equation}
D(s)\land P(s)\land R_+(s,s')\land\run{T}{s'}{o_t}
\ \Rightarrow\ \exists o_s.\ \run{S}{s}{o_s}\land o_t\preceq o_s.
\label{eq:conditional}
\end{equation}
The refused clause has $R_-$ and the fallback $F$, without requiring $P$.
When $F=S$, private-state transport still needs proof; syntactic reuse of the
source does not imply that it behaves correctly from $s'$.

\begin{theorem}[Local guarded refinement]
Let the guard and conditional certificates share a domain, premise, entry
relations, and branches. Every behavior of
$\operatorname{select}(g,T,F)$ from an entry satisfying $D$ is related to a
behavior of $S$ from that entry.
\end{theorem}
\begin{proof}
Apply the dispatch law to obtain the actual check result and checked state.
Guard soundness yields the appropriate entry relation and, on acceptance, $P$.
Apply the corresponding branch refinement clause.
\end{proof}
This is \code{guardify_refinement}. Source-to-target behavior preservation is
provided independently by \code{preservation_certificate} and
\code{guardify_preservation}, which also use check availability. Neither theorem
silently supplies the other direction or upgrades a finite big-step host to a
divergence theorem.

\subsection{Condition Libraries and Assumption Derivation}

The logical certificate chain distinguishes four links:
\begin{center}
\begin{tabular}{ll}
$C_{\mathrm{opt}}$&$A$ suffices for candidate correctness;\\
$C_{\mathrm{derive}}$&entry property $B$ implies the local/model obligations $A$;\\
$C_{\mathrm{guard}}$&the actual check safely establishes $B$ on acceptance;\\
$C_{\mathrm{host}}$&concrete guarded choice implements Eq.~\eqref{eq:dispatch}.
\end{tabular}
\end{center}
The kernel can consume the resulting certificate for $P=A$ without knowing how
$B$ was proposed. A transformation may supply a proof directly or submit an
untrusted candidate and witness to a verified checker.

Reusable condition services sit above the kernel. For example,
\code{sequence_readonly_conditions} composes a first check on $D$ with a second
check certified on $D\land P_1$. A refused first stage skips the second;
acceptance of both establishes $P_1\land P_2$. This stronger continuation domain
allows earlier checks to license later reads. Prefix-scan and probe-simplification
libraries expose further proof obligations rather than adding concrete memory
or optimizer semantics to the kernel.

\subsection{Local Reasoning and Complete Programs}

Language services establish local source/candidate execution bridges and the
frame that restores public observations. A domain implementation proves the
relation between the localized source and candidate models. Their composition
can discharge conditional equivalence, as implemented by
\code{localized_conditional_equivalence} in the read-only frontend.

Installation into a complete program is a separate language/IR obligation.
The generic \code{context_certificate} is a hook for a supplied lifting theorem:
it requires legal placement and admissibility of the actual replacement. It
does not prove that an arbitrary context preserves an arbitrary local relation.
The concrete Clight hosts establish control, private-resource, and progress
conditions and then compose with the CompCert backend.

Repeated rewrites are checked against their actual intermediate programs.
The read-only frontend's \code{certified_rewrite_sequence_equivalent} composes
finite sequences of installed equivalences. State relations, entry domains,
and placement evidence cannot simply be reused after an intervening rewrite
without showing that their requirements still hold.

\begin{table}[t]
\caption{Responsibilities at the proof boundary.}
\label{tab:responsibilities}
\centering\small
\begin{tabular}{>{\raggedright\arraybackslash}p{.19\textwidth}
  >{\raggedright\arraybackslash}p{.35\textwidth}
  >{\raggedright\arraybackslash}p{.36\textwidth}}
\toprule
Party&Supplies once or as a reusable library&Supplies for a transformation/site\\
\midrule
Kernel&Certificate interfaces and local composition&Consumes certificates for actual code/checks\\
Language/IR&Execution, primitive safety, dispatch, frames, installation laws&Site evidence satisfying host contracts\\
Domain implementer&Candidate checker, range/footprint and assumption services&Source/model bridge, sufficient premise, candidate and placement data\\
\bottomrule
\end{tabular}
\end{table}

\section{CompCert Case Study: Optimistic Affine Loops}
\label{sec:case-study}

The case study tests the interfaces against real Clight execution and memory.
PolCert supplies a capability reference for affine domains, candidate schedules,
and tiling; the integration uses GuardCert's own source/guard adapters together
with the existing vendored optimizer proof infrastructure. The concrete toolchain
is pinned to CompCert~3.18, Rocq~9.2, and Stdlib~9.2.

\subsection{Established Compiler Path}

The current complete path supports a loaded signed root bound with a constant
offset and canonical affine descendants. It captures the initial header word,
executes private numeric and memory checks, validates an untrusted candidate,
and retains the original source on refusal. Candidates include schedule
reordering and supported tiling/domain witnesses; successful checking binds the
candidate to the actual source description.

For this path, \code{compile_offset_affine_multi_regions_correct} proves that a
successful compiler result is related to the original Csem program by a
backward simulation into CompCert assembly semantics. The proof composes
SimplExpr, SimplLocals, the expression-region host, and the verified backend.
The region contracts and host progress obligations determine the source class;
the theorem is not a license to install arbitrary Clight fragments.

\subsection{A Nested Loaded-Bound Acceptance Target}

An adapted Figure~2 source from OLO~\cite{olo2017} supplies a concrete pending
acceptance target:
\begin{lstlisting}[language=C]
for (row = 0; row < shape[0] + 1; row++)
  for (column = 0; column < shape[1] + 1; column++)
    for (component = 0; component < 5; component++)
      output[(row * 16 + column) * 5 + component] =
          row + column + component;
\end{lstlisting}
This preserves the three-level control pattern and two repeatedly loaded
headers. It uses a flat fixed layout and an arithmetic store body; it is not
the full NAS BT kernel. The output may alias either header. A store can then
change the next source test, so the captured iteration domain need not describe
the original execution until preservation has been established.

The language services first capture the outer header. They capture the child
header only if the actual outer test is active. An empty outer loop therefore
requires neither a defined child counter nor a readable second header. The
source's modular \code{Int.add} computation is kept distinct from the raw memory
observation; no-wrap facts belong to the accepted modeling condition.

For an active prefix, the proof retains a witness in the memory actually reached
by the source. It transports access permissions back to the guard-entry memory,
not loaded data values. Each accepted body check must protect \emph{all} captured
observations before the next child test is justified. Completion of those checks
for the current row permits the next outer test. Only after this preservation
chain can the proof derive execution of the cached source and consume the
affine candidate certificate.

The ordered capture, indexed-load primitives, nested source-prefix services,
and two-cache transport are proved. A numeric probe can already run from
captured parameter words without assuming completion of the cached source.
Additional adapters now initialize private model bounds and transport execution
of a reached literal-bound subbody into the canonical affine model. They derive
that subdomain's access permissions at guard entry while retaining the reached
source memory. This bridge requires no stability assumption for the enclosing
loaded headers.

The domain adapter now consumes these permissions to execute a joint scan for
the reached constant-bound subbody. Observation descriptors carry actual address
expressions and load receipts, permitting comparisons against both
\code{shape} and \code{shape + 1} without another pointer temporary. Acceptance
establishes separation of every point write from every captured word, then
preservation for actual subbody executions with the same coordinate and pointer
view. This evidence discharges the inner-prefix preservation obligation.
A concrete fixture proves all five original stores and the actual scan: data
and headers overlapping are refused, while adjacent disjoint regions within one
allocation are accepted. It does not install an enclosing loop optimizer.

The subsequent adapter executes the inner short-circuit loop using source
prefix receipts for each column. Its acceptance supplies observation preservation
for every column and advances the current outer prefix. The source witness and
public guard state remain separate; mapped private cursors provide the scanned
coordinates. A refusal at the first body leaves the scan cursor at zero, and an
empty-child execution needs no body or output-pointer receipt. Fixed-entry prefix
services require the header law only at the captured entry. Numeric-domain,
word-view, and namespace facts remain explicit inputs until the factory produces
them from the original source and checked package.

The outer adapter now consumes this row producer in an actual short-circuit
loop. Complete acceptance supplies observation preservation for every active
row and column, discharging the preservation input to the existing two-cache
transport. The resulting cached source has the same exit temporaries and final
memory as the original source. A first-row refusal stops before the outer
increment. Separate empty-outer execution and zero-header transport services
need no child cache or body-pointer receipt. An unread child must not become a
premise merely to instantiate a two-observation prefix. These services do not
by themselves license a reordered candidate that reads the child first.
\pending{Assemble the canonical three-level affine model, the actual
capture/numeric producers for the typed metadata inputs, the original-source matcher/factory, and a new
whole-program compiler entry. The current adapted Figure~2 native probe retains
the original source and does not demonstrate an installed optimization.}

\subsection{Condition Expressiveness and Reuse}

The supported services cover machine-word ranges/no-wrap, byte-level memory
separation, permissions for typed accesses, and preservation of loaded values.
These are sufficient-condition libraries with declared syntax and semantic
domains. They do not decide arbitrary propositions or infer weakest conditions
for arbitrary program pairs. For a source-derived scan, the domain implementation
still proves access coverage and the source/model bridge. The language library
supplies the physical comparison and load/store preservation lemmas.

Condition construction must preserve partial evaluation safety. A mathematically
simpler condition can read a header or compare a pointer earlier than the source
permits. Reuse therefore requires a certificate for the reordered or simplified
check, not just Boolean equivalence on states where every expression is defined.
\pending{Select the final compact stability condition and compare its guard
work and useful acceptance against the existing scan on the same source. General
projection and arbitrary dependent preload expressions remain outside the
established implementation claims.}

\section{Related Work}
\label{sec:related}

The comparison separates condition discovery, safe executable checking,
conditional candidate proofs, and complete-program installation. Evidence for
one link does not imply all four, and names such as ``optimistic'' or ``peephole''
do not determine an interface's expressive power.

\paragraph{Optimistic loops and candidate conditioning.}
Doerfert et al.'s OLO~\cite{olo2017} organizes assumptions from polyhedral
modeling, derives and simplifies runtime conditions, and handles conditional
preloads and check arithmetic. These are direct algorithmic precedents for the
case study. Our comparison concerns a mechanical proof chain for a declared
Clight source class, not the invention of runtime loop versioning. OLO's
production compiler and performance results also set a practical standard that
our current proof endpoints alone do not meet.

COVE/cSTOKE~\cite{cove2015} accepts independently produced program candidates,
infers sufficient conditions using tests and analysis, and verifies conditional
equivalence. Its dynamic-check experiments use C instrumentation compiled and
joined with the optimized binary. This makes candidate conditioning a broader
motivation than loop optimization. GuardCert's current template-based adapters
do not provide COVE's general condition-inference capability. Our relevant
comparison is the proved connection between the actual generated check, its
safe evaluation domain, the candidate, and installation.

\paragraph{Verified speculation and recovery.}
CoreJIT~\cite{corejit2021} mechanically verifies speculative optimization and
deoptimization in an IR with explicit assumptions and synchronization points.
It separates optimization policy from verified execution mechanisms. This is a
precedent for verified dynamic assumptions, not merely an equivalence checker.
Our present entry-versioning contract retains an original fragment as fallback;
it does not reconstruct an active frame at an internal deoptimization point.
CoreJIT's 2021 native experiment lies outside that paper's proof boundary.
A later effectful-JIT development~\cite{fmjit2023} connects native generation to
the CompCert backend under its own scope. These two results should not be
combined into an unstated theorem about all speculative passes.

\paragraph{Local validation and whole-program lifting.}
Peek~\cite{peek2016} provides verified assembly peephole rules and a
liveness-based connection from local simulation to complete programs. Its
published framework imposes replacement/control restrictions and source
normalization obligations. These establish a substantial local-to-global
precedent. A comparison must identify the extra check-safety, private-state,
and source-prefix obligations of our loop example; moving the example to
Clight or adding a conditional is insufficient.

Chamois~\cite{chamois2023} validates oracle-proposed control-flow graphs using
block simulations and relational invariants, with integration into CompCert.
Its documented memory-copy expansion proves concrete CFG and private-resource
simulation~\cite{chamoisapi}. We have inspected those interfaces, but have not
instantiated its artifact with our loaded-bound example. The absence of such a
comparison is not evidence that its framework cannot express a guarded rewrite.

\paragraph{Comparison still required.}
\pending{For one fixed alias-sensitive loaded-bound source, record which
obligations and proof services OLO, CoreJIT, Chamois, Peek, and GuardCert supply.
Where artifacts are instantiated, report the additional source/model,
condition-safety, frame, and placement work. Distinguish unknown capabilities
from demonstrated restrictions before making a novelty or author-effort claim.}

\section{Evaluation Plan and Established Evidence}
\label{sec:evaluation}

The evaluation distinguishes semantic proof, compilation, execution validation,
and cost. A successful native test is not a universal correctness proof, and
proof counts are not a measure of practical optimization value.

\subsection{Current Artifact Evidence}

For the loaded-root-with-offset compiler, the existing report records six
configurations and 708 full assembly calls. Separate path probes record 236
Clight calls and 21 probes of unmodified assembly. These cover candidate
acceptance and fallback, changing bounds through aliases, empty paths,
multi-array bodies, and repeated rewrites. The reports are bound to their
compiler/source digests; these are prior execution results, not a new matrix run
for this manuscript.

The earlier Figure~2 coverage report records sixteen unchanged-source calls;
it remains a fixed record, not an optimization-success count. The new frontend
closure covers 20 endpoints: 17 language results, one domain result, and two
fixtures, with 569 required dependencies and 1,104 source digests. It binds the
parent nested compiler closure, whose 25 endpoints connect original-source
checking, canonical execution, alias/candidate validation, kernel preservation,
and installation. The kernel is unchanged; the new compiler retains the same
42-global assumption set, with no additional global axiom.

The new extracted frontend compiler was run on the same adapted C source.
Six configurations---disabled, identity, interchange, $2\times3$ tiling, wrong
reindexing, and invalid domain---give 48 complete assembly calls. A further 27
calls use a harness with one additional input and unchanged kernel/context
prefix. All 75 calls compare all 2,048 output words, both headers, and public
counter exits against a GCC wrapping-arithmetic reference and an independent
source-word model. Separate instrumented Clight runs give sixteen dispatch
observations: two fast and six fallback calls for each reordered configuration.
These Clight observations are not assembly-path evidence.

Six debugger probes observe unmodified assembly. Three watch selected output
locations and distinguish their write order for one independent-array input
(Table~\ref{tab:nested-order}); the
other three check first-header or second-header-region alias fallback. The
second-header-region probe excludes a write that the cached iteration domain
would have made. Public exits and full outputs match source behavior. Probe
runs are separate from the 75-call assembly matrix. Exact commands, source and
compiler bindings, report hashes, and historical stages are in the evidence map.

\begin{table}[t]
\centering
\caption{Write order at three watched output locations in assembly for one
independent-array input. All three runs finish with public counters (3, 4, 5).}
\label{tab:nested-order}
\begin{tabular}{llr}
\toprule
Configuration & Output word indices & Function bytes \\
\midrule
Disabled & 15, 80, 160 & 161 \\
Interchange & 80, 160, 15 & 564 \\
2-by-3 tiling & 80, 15, 160 & 693 \\
\bottomrule
\end{tabular}
\end{table}
The identity configuration uses 528 function bytes. Code size is an observation
of the current generated programs, not a timing or profitability result.

\subsection{Experiments Needed for the Final Paper}

\paragraph{Semantic coverage.}
Classify accepted candidates, runtime refusals, and static matcher/validator
refusals. Use the same source for independent arrays, header/data aliasing,
empty outer and inner loops, overflowing offset expressions, limited readable
allocations, invalid candidate domains, and conflicting dependencies. Verify
complete arrays and public loop exits, and inspect executed store order to
distinguish a real fast path from unchanged-source compilation.

\paragraph{Condition construction.}
Compare the scan and compact sufficient conditions on a common source class.
Report accepted input domains and refused cases. Count captures, pointer
comparisons, numeric operations, and iterations separately. Existing work
counts show that loop-shaped guard code can still perform expensive repeated
comparisons; code-size reduction does not establish cheaper checking.

\paragraph{Runtime and compilation cost.}
Use the same CompCert version, target, and backend configuration for the original
and transformed programs. Measure check-only work where separable and total
execution cost, including refusal and empty paths. Report machine code size,
compile time, measurement environment, raw batches, and dispersion. Avoid
concurrent proof/build work during timing. Slowdowns and narrow acceptance are
results to explain rather than cases to remove.

\paragraph{Proof reuse.}
For multiple transformations, identify the shared kernel, language lemmas, and
domain services. List the site-specific assumptions and source/model bridges
still supplied by the author. Compare a common example against the nearest
existing interfaces before claiming reduced proof burden.

\pending{No timing, speedup, final acceptance fraction, or comparative author
effort has been measured for this manuscript. Extend this frontend's coverage to
multi-array read/write bodies, dependency refusals, same-allocation slices,
nonzero entry counters, and repeated/control contexts; validate a fresh-build
bootstrap. Measure compact conditions and author obligations on common examples.}

\section{Conclusion and Remaining Work}

GuardCert's implemented kernel composes executable-check and transformation
certificates without fixing a source language or optimizer. Its CompCert instance
makes check safety, private-state transport, and installation concrete for
supported loop transformations. The pending nested loaded-bound path tests
whether these services remove the circularity between source execution and
guard justification. Closing that path, constructing affordable sufficient
conditions, and measuring reuse and cost will determine the final case-study
and contribution claims.

\bibliographystyle{splncs04}
\bibliography{references}
\end{document}
